\begin{mvtable}
\hfill{\small 续表\par}\vspace{2pt}
\begin{longtable}{cccccc}\toprule
序号 & 电压($X_1$) & 电流($X_2$) & 进给速度($X_3$) & 气流($X_4$)\\
\midrule
\endfirsthead
\multicolumn{5}{c}{\zihao{6}续表}\\
\toprule
序号 & 电压($X_1$) & 电流($X_2$) & 进给速度($X_3$) & 气流($X_4$)\\
\midrule
\endhead
\bottomrule
\endlastfoot

9 & 22.5 & 277 & 289.0 & 52.0\\
10 & 22.5 & 278 & 289.0 & 52.0\\
11 & 22.3 & 269 & 287.0 & 54.0\\
12 & 21.8 & 274 & 287.6 & 52.0\\
13 & 22.3 & 270 & 288.4 & 51.0\\
14 & 22.2 & 273 & 290.2 & 51.3\\
15 & 22.1 & 274 & 286.0 & 51.0\\
16 & 22.1 & 277 & 287.0 & 52.0\\
17 & 21.8 & 277 & 287.0 & 51.0\\
18 & 22.6 & 276 & 290.0 & 51.0\\
19 & 22.3 & 278 & 287.0 & 51.7\\
20 & 23.0 & 266 & 289.1 & 51.0\\
21 & 22.9 & 271 & 288.3 & 51.0\\
22 & 21.3 & 274 & 289.0 & 52.0\\
23 & 21.8 & 280 & 290.0 & 52.0\\
24 & 22.0 & 268 & 288.3 & 51.0\\
25 & 22.8 & 269 & 288.7 & 52.0\\
26 & 22.0 & 264 & 290.0 & 51.0\\
27 & 22.5 & 273 & 288.6 & 52.0\\
28 & 22.2 & 269 & 288.2 & 52.0\\
29 & 22.6 & 273 & 286.0 & 52.0\\
30 & 21.7 & 283 & 290.0 & 52.7\\
31 & 21.9 & 273 & 288.7 & 55.3\\
32 & 22.3 & 264 & 287.0 & 52.0\\
33 & 22.2 & 263 & 288.0 & 52.0\\
34 & 22.3 & 266 & 288.6 & 51.7\\
35 & 22.0 & 263 & 288.0 & 51.7\\
36 & 22.8 & 272 & 289.0 & 52.3\\
37 & 22.0 & 277 & 287.7 & 53.3\\
38 & 22.7 & 272 & 289.0 & 52.0\\
39 & 22.6 & 274 & 287.2 & 52.7\\
40 & 22.7 & 270 & 290.0 & 51.0\\
\end{longtable}
\end{mvtable}

资料来源:数据经 Mark Abbotoy 许可使用.

对于大多数变量来说,正态的假设是合理的,但是我们对气流数据取了自然对数. 另外,对于每个变量的连续观测,都没有明显的序列相关性.

\hyperref[fig:5.9]{图 5.9} 是四个焊接变量的 $T^{2}$ 控制图. 虚线是 $95\%$ 的限制,实线是 $99\%$ 的限制. 采用 $99\%$ 的限制,没有点会在控制区域外. 而对于 $95\%$ 的限制,序号为 31 的点在控制区域外.

对于两个变量,椭圆质量控制图(椭圆控制图)告诉我们什么呢? 大多数变量在控制图之
